\section{The fixed prior and the main guarantee}
\label{sec:sharp-model}
Put $D=d+1$ and $\bar x=(1,x)$ for $x\in[-1,1]^d$. The learner has one
prescribed hidden width $N$ and functions
\[
 f_{a,u}(x)=\sum_{i=1}^N a_i\phi(u_i\cdot\bar x),
\]
where $\phi$ is $1$-Lipschitz and $|\phi|\le1$. Tanh is the main example.
There is no final output link. The prior is supported on $\|a\|_1\le1$,
so every raw output belongs to $[-1,1]$.

Fix $p>2$ and $0<c_0\le\min\{1,(p-1)/8\}$. For incoming feature
coordinates, use the symmetric finite-variation L\'evy measure
\[
 \nu(dz)=\frac{c_0\,dz}{|z|(1+|z|)^p}.
\]
Let $\vartheta_F=\mathcal L(L_{1/F})$ be its zero-drift Poisson
increment law. All incoming coordinates independently have law
$\vartheta_D$. Before the readout norm conditioning, the readout
coordinates independently have the signed beta-prime law $\lambda_N$:
its sign is fair and its independent magnitude has density
\[
 q_a(t)=\frac{t^{a-1}(1+t)^{-a-p}}{\mathrm B(a,p)},\qquad
 t>0,\qquad a=c_0/N.
\]
Write $\pi_N^0$ for this product parameter prior and define
\[
 A_N=\{\|a\|_1\le1\},\qquad p_N=\pi_N^0(A_N),\qquad
 \pi_N=\pi_N^0(\cdot\mid A_N).
\]
Because $\mathbb E_{\lambda_N}|a_i|=c_0/[N(p-1)]$,
$p_N\ge1-c_0/(p-1)\ge7/8$. Incoming features remain independent of the
conditioning. The readout coordinates become dependent, but remain
nonzero almost surely. Their unconditioned variance is
$a(a+1)/[(p-1)(p-2)]$, of order $1/N$. These beta-prime readout laws
are not asserted to be L\'evy increments. All scalar constants are
fixed independently of sample size and architecture; only fan-in
changes their prescribed coordinate laws.

A teacher has the form
\[
 g(x)=\sum_{j=1}^C b_j\phi(v_j\cdot\bar x),\qquad
 \sum_j|b_j|\le1,\quad \|v_j\|_\infty\le B,\quad B\ge1.
\]
The incoming cap is removed for tanh in
Section~\ref{sec:l1-compactification}. The readout variation restriction
remains part of both the learner and teacher models.

Let $(X_i,Y_i)$ be IID, with $X_i\in[-1,1]^d$,
$f_0(x)=\mathbb E[Y\mid X=x]\in[-1,1]$, and conditionally centered
$\sigma^2$-sub-Gaussian noise. Put $s=1+\sigma^2$ and
$R(f)=\mathbb E(f(X)-Y)^2$. Fix $T=2048s$ and use the ordinary posterior
\[
 dQ_S(w)=Z_S^{-1}\exp\left\{-\frac1T
                  \sum_{i=1}^n(f_w(X_i)-Y_i)^2\right\}d\pi_N(w),
 \qquad \bar f_S=\int f_w\,dQ_S(w).
\]
Every observation enters this likelihood, whose Gaussian variance is
$T/2=1024s$. Architecture and prior are fixed before observing data.

\begin{theorem}[Sharp high-probability all-data oracle]
\label{thm:l1-main}
For every deterministic teacher as above, $n\ge2$, $0<\delta<1$, and
every prescribed finite $N\ge8C$, with probability at least $1-\delta$,
\begin{align}
 \int R(f_w)\,dQ_S(w)
 &\le R(g)+K\frac C{N^2}\notag\\
 &\quad+\frac{Ks}{n}\left[(1+CD)\log\{eD(B+8)(C+1)n\}+\log\frac4\delta\right].
 \label{eq:sharp-main-bound}
\end{align}
Here $K$ depends only on $c_0,p$. The final posterior mean satisfies the
same bound. If $N\ge\max\{8C,\sqrt n\}$, the $C/N^2$ term is absorbed
into the statistical remainder.
\end{theorem}

The theorem makes no computational-efficiency claim. There is no
selection over architectures, holdout, temperature selection, or
averaging over sample sizes. Its comparison law and all precision
choices below are proof devices, not changes to the learner.

\section{Scalar laws and feature placement}
The L\'evy measure has infinite total mass but finite first and second
absolute moments. Its Poisson sum is absolutely convergent. Each
$\vartheta_F$ is absolutely continuous: for every cutoff $r>0$, the
large-jump sum has a density conditional on at least one jump, and the
probability of no such jump tends to zero as $r\downarrow0$. Full support
follows by prescribing one jump near any desired value and requiring
the independent small-jump remainder to be small. Thus the incoming
prior has full support and continuous coordinates.

For $F\ge1$, $|b|\le B$, and $0<r\le1$, there is a law
$q_{F,b,r}\ll\vartheta_F$ such that
\begin{equation}
 |W-b|\le12r,\qquad
 \mathrm{KL}(q_{F,b,r}\|\vartheta_F)
 \le\log F+C_0\log\frac{C_0(B+8)}r.
 \label{eq:l1-scalar-placement}
\end{equation}
When $|b|>4r$, the stronger support bound $|W-b|\le r$ holds.
For completeness, split the Poisson process at $r/\kappa$, with fixed
large $\kappa$. Prescribe exactly one large jump, uniform on
$[b-r/2,b+r/2]$ if $|b|>4r$. Otherwise use a mixture of uniforms on
$[5r,7r]$ and $[-7r,-5r]$ with mean $b$. Retain the original small-jump
process conditioned on its sum having magnitude at most $r/2$.
Its failure probability is $O(1/F)$ by its variance, so this conditioning
costs $O(1/F)$. The prescribed jump has density at most $1/r$, and the
original Levy density on its support is at least a constant times
$(B+8)^{-p-1}$. The entropy on Poisson space is
\[
 \int q(z)\log\frac{Fq(z)}{\nu(z)}\,dz
 +\frac{\nu\{|z|>r/\kappa\}}F
 -\log\Pr\{|S_{r/\kappa}|\le r/2\},
\]
which gives \eqref{eq:l1-scalar-placement}; summing the jumps decreases
KL. The support bounds follow directly from the displayed jump intervals.
The small-$b$ mixture can be chosen with mean $b$ because $|b|\le4r<6r$.

In particular, every feature with $\|v\|_\infty\le B$ has a comparison
law $\zeta_{v,\epsilon}$ relative to its native product incoming prior
such that
\begin{equation}
 \sup_x|\phi(U\cdot\bar x)-\phi(v\cdot\bar x)|\le\epsilon,
 \qquad
 \mathrm{KL}(\zeta_{v,\epsilon}\|\vartheta_D^{\otimes D})
 \le C_0D\log\frac{C_0D(B+8)}\epsilon.
 \label{eq:l1-feature-placement}
\end{equation}
Apply \eqref{eq:l1-scalar-placement} independently with
$r=\epsilon/(12D)$ and use Lipschitzness.

\subsection{Scaling, placement, and merging of scalar coordinates}
For $0<s<1$, direct density division gives
\begin{equation}
 \mathcal L(sW)\le s^{-a}\lambda_N,
 \qquad W\sim\lambda_N.
 \label{eq:merge-beta-scaling}
\end{equation}
Indeed the density ratio is
$s^{-a}[(1+|w|)/(1+|w|/s)]^{a+p}\le s^{-a}$.
For $0<a\le c_0\le1$, the quantity $a\mathrm B(a,p)$ is bounded above
and below by fixed positive constants. Also
\[
 t^{a-1}(1+t)^{-a-p}
 =(t/(1+t))^{a-1}(1+t)^{-p-1}\ge(1+t)^{-p-1}.
\]
Therefore a uniform distribution on $[b-r,b+r]$, with $|b|\le1$ and
$0<r\le1$, has density relative to $\lambda_N$ at most $C_pN/r$.
It has mean $b$ and lies within $r$ of its center. The incoming-feature
placement law from \eqref{eq:l1-feature-placement} is unchanged.

The positive density $q_a$ is decreasing because $a\le1$. Hence
\begin{equation}
 (q_a*q_a)(t)
 \le2q_a(t/2)\int_0^{t/2}q_a(u)\,du
 \le2^{p+2}q_a(t).
 \label{eq:merge-beta-convolution}
\end{equation}
The last inequality uses
$q_a(t/2)/q_a(t)\le2^{p+1}$. Thus the symmetric law whose magnitude is
the sum of two independent native magnitudes has density at most
$2^{p+2}$ times the native signed law.

For any fixed parameter pair $(a_i,u_i),(a_j,u_j)$, put
$L=|a_i|+|a_j|$. Choose feature $u_i$ with probability $|a_i|/L$ and
feature $u_j$ otherwise. In the first slot place the chosen feature
and coefficient $\operatorname{sign}(a_{\rm chosen})L$; erase the second
slot. At $L=0$ use any choice and coefficient zero. This operation
preserves the pair's $\ell^1$ norm pointwise and its function in conditional
expectation. At every $x$ its conditional variance is
\begin{equation}
 |a_i a_j|\{\operatorname{sign}(a_i)\phi(u_i\cdot\bar x)
       -\operatorname{sign}(a_j)\phi(u_j\cdot\bar x)\}^2
 \le4|a_i a_j|.
 \label{eq:merge-pair-variance}
\end{equation}
Under the unconditioned product reference, the copied feature retains
its native law and is independent of the new coefficient: the choice
depends only on readout magnitudes, the two features are iid, and the
chosen sign is fair independently of those magnitudes. Thus a fixed
pair merge has precisely the scalar domination
\eqref{eq:merge-beta-convolution}; no orientation factor is required
because the retained position is always the first slot.

\section{A finite-width midpoint transfer}
\begin{lemma}[Unbiased finite-width midpoint transfer]
\label{lem:merge-transfer}
Let $N\ge8C$, $0<\epsilon\le1$, and let $g$ obey the same teacher
width, incoming coefficient and readout variation restrictions as stated above.
For every finite-KL $Q\ll\pi_N$ there is a law
$Q_g\ll\pi_N$ coupled with source and output functions $F,H$, such that
\begin{align}
 \mathrm{KL}(Q_g\|\pi_N)
 &\le\mathrm{KL}(Q\|\pi_N)+J(C,B,\epsilon),
 \label{eq:merge-KL}\\
 \left|\mathbb E[H(x)-(F(x)+g(x))/2\mid F\text{ parameters}]\right|
 &\le\epsilon,
 \label{eq:merge-mean}\\
 \mathbb E[|H(x)-(F(x)+g(x))/2|^2\mid F\text{ parameters}]
 &\le C_p C/N^2+C_p\epsilon^2.
 \label{eq:merge-second}
\end{align}
The last two inequalities hold uniformly in $x$, and one may take
$J=C_p[1+CD\log\{eC_pD(B+8)(C+1)/\epsilon\}]$.
The same conditional bounds hold after fixing the selected slot indices.
\end{lemma}
\begin{proof}
Set $r=\epsilon/[64(C+1)]$, $\gamma=4Cr$, and retain teacher coefficients
with magnitude greater than $16r$, leaving $k\le C$ terms. At every old
readout vector of $\ell^1$ norm at most one there are at least $N/2$ slots
with $|a_i|\le2/N$. Choose an ordered tuple of $2k$ distinct indices
uniformly among such slots and group consecutive entries into $k$ pairs.
Relative to a uniform tuple among all $(N)_{2k}$ possibilities, this
conditional choice costs at most $2k\log4$ in KL: at each draw there
are at least $N/2-2C\ge N/4$ permitted remaining indices.

Independently merge each pair as above, always retaining its first
position. The resulting function $F^{\rm m}$ has conditional mean $F$,
and, because the merge choices are independent and each pair has
$|a_i|,|a_j|\le2/N$,
\[
 \mathbb E[|F^{\rm m}(x)-F(x)|^2\mid W,I]\le16k/N^2.
\]
Its readout $\ell^1$ norm equals the old one. Halve all remaining readout
coefficients and use the freed second position in pair $j$ for a feature
from \eqref{eq:l1-feature-placement} and a uniform coefficient within
$r$ of $(1-\gamma)b_j/2$. The new readout norm is at most
$\frac12+\frac12(1-\gamma)+Cr\le1$, and every inserted
magnitude exceeds $r$. Relative to $g/2$, discarded teacher terms
cost at most $8Cr$, shrinkage costs at most $2Cr$, readout perturbations
at most $Cr$, and feature perturbations at most $\epsilon/2$.
Thus the inserted teacher function differs uniformly from $g/2$ by at
most $11Cr+\epsilon/2\le\epsilon$. This insertion error is independent of
the merge choices. Its conditional mean is bounded by $\epsilon$ and
its second moment by $\epsilon^2$. Halving the merging error and using
its zero conditional mean proves \eqref{eq:merge-mean} and
\eqref{eq:merge-second}, with a numerical enlargement of $C_p$.

Here is an explicit density proof of the KL bound. Augment the source
law with the selected tuple. Relative to $\pi_N$ times a uniform
tuple, its KL is at most $\mathrm{KL}(Q\|\pi_N)+2k\log4$.
Apply the same randomized merging and insertion kernel to both laws;
data processing bounds the output KL relative to the transformed
reference by this number.

For a fixed tuple, the unconditioned reference after merging has native
independent features, native readout coordinates outside the retained
positions, and $k$ merged readout laws dominated by $2^{p+2}\lambda_N$.
The merged magnitude exactly preserves the old $\ell^1$ norm. Consequently
conditioning the old prior on its norm can be bounded by the same
unconditioned reference divided by its old normalizer. Halving all
$N-k$ retained coefficients costs at most $2^{c_0}$ by
\eqref{eq:merge-beta-scaling}. Inserting the $k$ teacher coefficients
costs at most $(C_pN/r)^k$, while inserting their features costs at most
\[
 \exp\{C_p kD\log(C_pD(B+8)/\epsilon)\}.
\]
For this last density claim, the feature placement construction also
has an L-infinity density bound of the displayed order: its prescribed
one-jump Poisson density and its small-jump conditioning density are
pointwise bounded by the same factors used in the entropy estimate.
Taking a pushforward preserves an upper density bound.

The output is supported on the same $\ell^1$ ball, so dividing by the target
conditioned prior cancels the old conditioning normalizer exactly.
At a realized output, at most $1/r$ slots have readout magnitude greater
than $r$. Thus the $k$ inserted positions have at most $r^{-k}$ possible
ordered choices. The $k$ retained first positions have at most $N^k$
choices. Averaging over all $(N)_{2k}$ tuples bounds the transformed
reference density relative to $\pi_N$ by
\[
 2^{c_0}\,2^{(p+2)k}
 \frac{N^{2k}}{(N)_{2k}}\,r^{-k}
 \exp\{C_pkD\log(C_pD(B+8)(C+1)/\epsilon)\}.
\]
Since $N\ge8C\ge8k$, $N^{2k}/(N)_{2k}\le(4/3)^{2k}$.
The logarithm of this density bound, plus the selection cost, is at
most $J(C,B,\epsilon)$. Combining it with data processing proves
\eqref{eq:merge-KL} without multiplying the source KL by a constant.
The argument does not assume independence of the source posterior.
\end{proof}

% Uniform-width direct localization and the bounded-variation uniform bound.
\section{Direct localization at every realizable width}
\label{sec:direct-localization}
We use the signed beta-prime readout prior and the same incoming
feature laws and $\ell^1$ conditioning fixed in Section~\ref{sec:sharp-model}.

\begin{lemma}[Small beta-prime backgrounds]
\label{lem:beta-background}
There are constants $r_0,K>0$ depending only on $c_0,p$ such that, for
$a=c_0/N$ and $0<r\le r_0$, a native readout magnitude $V$ satisfies
\[
 -\log\Pr(V\le r)\le Ka\log(K/r),\qquad
 \mathbb E[V\mid V\le r]\le2ar.
\]
\end{lemma}
\begin{proof}
Put $p_r=\Pr(V\le r)$. Since $(1+t)^{-a-p}\ge1-(a+p)t$,
\[
 p_r\ge\frac{r^a}{a\mathrm B(a,p)}
       \left[1-\frac{(a+p)ar}{a+1}\right].
\]
Choose fixed $r_0$ so the bracket is at least $1/2$. The function
$\log[a\mathrm B(a,p)]$ extends smoothly to $a=0$ with value zero;
its derivative is bounded on $[0,c_0]$. Thus its absolute value is at
most $Ka$, while the negative logarithm of the bracket is at most
$Kar$. This proves the first bound. For the conditional mean, the
numerator is at most $r^{a+1}/[(a+1)\mathrm B(a,p)]$. Divide by the
preceding lower bound and use its bracket at least $1/2$.
\end{proof}

\begin{lemma}[Width-uniform sup-norm comparison]
\label{lem:direct-comparison}
For every $N\ge C$, $0<\epsilon\le1$, and every teacher $g$ of width
$C$, incoming cap $B\ge1$, and readout variation at most one, there is
a data-independent law $\rho_{g,\epsilon}\ll\pi_N$ such that
\begin{align}
 \rho_{g,\epsilon}\{\|f-g\|_\infty\le\epsilon\}&=1,\notag\\
 \mathrm{KL}(\rho_{g,\epsilon}\|\pi_N)
 &\le K(1+CD)\log\frac{eKD(B+8)(C+1)}\epsilon.
 \label{eq:direct-comparison}
\end{align}
The constant is independent of $N$.
\end{lemma}
\begin{proof}
Choose a sufficiently large fixed $K_0$ and set
$r=\epsilon/[K_0(C+1)]\le r_0$. Discard teacher readout coefficients
with magnitude at most $16r$, leaving $k\le C$ terms. Center the
retained coefficient laws at $(1-\epsilon/2)b_j$, with uniform radius
$r$. The centers have magnitude at least $8r$, so these inserted
coefficients have magnitude greater than $r$. Approximate each retained
feature uniformly to $\epsilon/8$ using the incoming placement law.

For any fixed selection of $k$ distinct slots, condition every other
native readout coordinate to have magnitude at most $r$, and then
condition their total $\ell^1$ norm to be at most $\epsilon/8$. By
Lemma~\ref{lem:beta-background}, their conditional expected total norm
is at most $2c_0r$. For large enough $K_0$, the second conditioning
therefore has probability at least $1/2$. The joint background law has
density relative to its native product reference at most
\[
 2p_r^{-(N-k)}\le\exp\{K\log(K/r)\},
\]
uniformly in $N$.

Every resulting readout has norm at most
$1-\epsilon/2+Cr+\epsilon/8\le1$. The uniform function error consists
of at most $16Cr$ from discarded coefficients, $\epsilon/2$ from
shrinkage, $Cr$ from coefficient perturbations, $\epsilon/8$ from
feature approximation, and $\epsilon/8$ from the background. For
$K_0\ge128$, their sum is at most $\epsilon$.

It remains to cancel the cost of selecting the slots. When $k\ge1$,
partition the $N$ slots into $k$ deterministic groups, each of size at
least $\lfloor N/k\rfloor\ge N/(2k)$, and choose one slot uniformly
from each group for its designated teacher term. For a fixed choice,
the readout insertion densities cost at most $(KN/r)^k$, and the
feature insertion densities cost at most
$\exp\{KkD\log(KD(B+8)/\epsilon)\}$. All other readouts have magnitude
at most $r$, whereas each inserted readout has magnitude greater than
$r$. The realized weights therefore identify the chosen position in
each group. In the mixture density, at most one placement contributes;
its uniform mixing weight cancels $N^k$ up to $(2k)^k$.

Combining these density bounds gives an upper density ratio of
$\exp\{K(1+CD)\log(eKD(B+8)(C+1)/\epsilon)\}$ relative to the
unconditioned product prior. The comparison is supported on the l1
ball, so converting the denominator to the conditioned prior multiplies
this ratio by $p_N\le1$. Its KL is bounded by the logarithm of the
displayed ratio. For $k=0$, use only the background law; the same bounds
hold without groups or insertion costs.
\end{proof}


\section{From midpoint transfer to a sharp population oracle}
\label{sec:sharp-statistics}
\begin{proof}[Proof of Theorem~\ref{thm:l1-main}]
Write $\beta=1/T$ and $\Delta_f=(f-Y)^2-(g-Y)^2$. Use
Lemma~\ref{lem:merge-transfer} with $\epsilon=n^{-2}$ and
$J=J(C,B,n^{-2})$. Conditional on the source parameters, expand squared
loss about the midpoint $M=(F+g)/2$. The conditional mean error and
second-moment bounds in that lemma give
\begin{align*}
 \mathbb E_{Q_g}\widehat R(H)
 &\le\tfrac12\mathbb E_Q\widehat R(F)+\tfrac12\widehat R(g)
       -\tfrac14\mathbb E_QP_n(F-g)^2+\varepsilon_S,\\
 \varepsilon_S&\le K C/N^2+K(1+P_n|Y|)n^{-2}.
\end{align*}
Gibbs variational optimality, together with the additive KL transfer,
therefore implies
\begin{equation}
 \mathbb E_QP_n\Delta_F
 \le-\tfrac12\mathbb E_QP_n(F-g)^2
            +2\varepsilon_S+2TJ/n.
 \label{eq:sharp-empirical-offset}
\end{equation}

We prove the needed relative-loss exponential inequality explicitly.
For fixed bounded $f,g$, write $d_f=f-g$, $D_f=d_f^2$, and
$B_f=d_f(f+g-2f_0)$. Then
$\Delta_f=B_f-2\xi d_f$, $B_f^2\le16D_f$, and $0\le D_f\le4$.
For $\kappa_*=1/8$ and $0<\lambda\le1/(1024s)$,
\begin{equation}
 \mathbb E\exp\{\lambda[P\Delta_f-\Delta_f-\kappa_*D_f]\}\le1.
 \label{eq:sharp-relative-mgf}
\end{equation}
Indeed, conditional integration of the noise leaves
$W=-B_f-\kappa_*D_f+2\lambda\sigma^2D_f$ inside $e^{\lambda W}$.
At these constants, $|W|\le6|d_f|$ and $|\lambda W|\le1$.
Since $e^z\le1+z+z^2$ for $|z|\le1$,
\[
 \mathbb E e^{\lambda W}
 \le1-\lambda PB_f-\lambda\kappa_*PD_f
       +\lambda^2(2\sigma^2+36)PD_f
 \le e^{-\lambda PB_f}.
\]
This proves \eqref{eq:sharp-relative-mgf}. Take $\lambda=2\beta$,
use independence, integrate over the prior, and apply Markov. With
probability at least $1-\delta/2$,
\begin{equation}
 \int e^{2\beta n[P\Delta_f-P_n\Delta_f-\kappa_*P_nD_f]}
                   d\pi_N(f)\le2/\delta.
 \label{eq:sharp-mgf-event}
\end{equation}

Put $L=\log(2/\delta)$. The conditional sub-Gaussian inequality
$\mathbb E e^{t|\xi|}\le2e^{\sigma^2t^2/2}$ and independence give, with
probability at least $1-\delta/2$,
\[
 P_n|Y|\le1+\sigma\sqrt{2(\log2+L/n)}.
\]
If $L>n$, the theorem follows from $R(f)-R(g)\le4$ by enlarging $K$.
Otherwise this event implies $1+P_n|Y|\le K\sqrt s$.

Use the direct comparison law of Lemma~\ref{lem:direct-comparison}
at accuracy $n^{-2}$. Its uniform approximation gives
\[
 \mathbb E_\rho P_n\Delta_f\le2(1+P_n|Y|)n^{-2}
                         \le K\sqrt s/n^2.
\]
Its KL is at most
$J_{\rm loc}=K(1+CD)\log\{eKD(B+8)(C+1)n\}$.
For $Z_g=\int e^{-\beta nP_n\Delta_f}\,d\pi_N(f)$, the variational
inequality therefore yields
\begin{equation}
 -\log Z_g\le J_{\rm loc}+K=:L_g.
 \label{eq:sharp-partition}
\end{equation}
This comparison law does not depend on the observed sample.

Rewrite \eqref{eq:sharp-mgf-event} using
$dQ_S/d\pi_N=e^{-\beta nP_n\Delta_f}/Z_g$, and apply Jensen.
Together with \eqref{eq:sharp-partition}, this gives
\[
 \mathbb E_QP\Delta_f
 \le\tfrac12\mathbb E_QP_n\Delta_f+\kappa_*\mathbb E_QP_nD_f
            +\frac T{2n}\{L_g+\log(2/\delta)\}.
\]
Insert \eqref{eq:sharp-empirical-offset}. The coefficient of the
empirical distance is $-1/4+\kappa_*=-1/8$, so this term can be dropped.
Consequently
\[
 \int R(f)\,dQ_S\le R(g)+\varepsilon_S
        +\frac Tn\{J+L_g/2+\tfrac12\log(2/\delta)\}.
\]
Substitute the bounds on $\varepsilon_S,L_g$ and the displayed
expression for $J(C,B,n^{-2})$. This is
\eqref{eq:sharp-main-bound}. Jensen proves the posterior-mean claim.
\end{proof}

\paragraph{Simultaneous comparison classes.}
For deterministic countably many teacher classes with uniform resources
$(C_j,B_j)$ and $8C_j\le N$, choose weights $w_j>0$ summing to at most
one. In each class choose a deterministic teacher within $1/n$ of its
population-risk infimum, apply the theorem with confidence $\delta w_j$,
and union bound. The same single posterior then obeys the infimum of
\eqref{eq:sharp-main-bound} over those classes, with
$\log(4/\delta)$ replaced by $\log(4/(\delta w_j))$ and an absorbed
$1/n$ error. These confidence weights are proof devices and are not
supplied to the learner. Dyadic incoming bounds compare arbitrary finite
teacher coefficients with their corresponding resources.

